%%%%%%%%%%%%%%%%%%%%%%%%%%%%%%%%%%%%%%%%%%%%%%%%%%%%%%%%%%%%%%%%%%%%%%%%%%%%%%%%
% METHODOLOGY: Hybrid Classical-Quantum Transfer Learning Benchmark
%%%%%%%%%%%%%%%%%%%%%%%%%%%%%%%%%%%%%%%%%%%%%%%%%%%%%%%%%%%%%%%%%%%%%%%%%%%%%%%%

\section{Methodology}
\label{sec:methodology}

Our experimental framework is designed to provide a rigorous, statistically robust comparison between classical and quantum transfer learning strategies in the presence of realistic hardware noise. This section details the architectural choices, baseline justifications, and statistical protocols required to address various reviewer concerns regarding scientific fairness and hardware reproducibility.

\subsection{Hybrid Architecture and Feature Extraction}
We employ a \textit{Hybrid Classical-Quantum Transfer Learning} approach. Pre-trained convolutional neural networks (ResNet-18, MobileNet-V2, EfficientNet-B0, and RegNet-400MF) serve as frozen feature extractors ($\phi$). Given an input image $\mathbf{x}$, the extracted feature vector $\mathbf{v} = \phi(\mathbf{x})$ is mapped to a low-dimensional Hilbert space via a Variational Quantum Circuit (VQC) head.

The VQC architecture ($h_\theta$) consists of:
\begin{enumerate}
    \item \textbf{Classical Projection:} A linear layer $W \in \mathbb{R}^{d \times n}$ where $d$ is the feature dimension and $n=4$ is the number of qubits. This reduces the search space before quantum encoding.
    \item \textbf{Angle Encoding:} Features are encoded into rotation angles $\alpha_i = \arctan(v_i) \cdot \pi/2$ to prevent saturation in the $RY$ gates.
    \item \textbf{Variational Layers:} $L$ layers of \textit{Strongly Entangling Layers}, consisting of three Euler rotations ($RZ, RY, RZ$) per qubit followed by a circular CNOT entanglement chain.
\end{enumerate}

\subsection{Scientific Fairness: Parameter-Matched Baselines}
A common critique in hybrid quantum literature is the use of oversized classical baselines (e.g., standard MLPs with thousands of parameters) compared to small quantum circuits. To ensure a fair comparison, we introduce a \textbf{Parameter-Matched MLP (MLP-SM)}. While a standard MLP-LG (128-64 units) contains over 70,000 parameters, our MLP-SM is constrained to match the total trainable weights of the VQC head exactly ($\approx 2,094$ parameters). This allows for an isolation of the "Quantum Advantage" property versus simple parameter scaling.

\subsection{Realistic Hardware Noise Modeling}
To address concerns regarding hardware fidelity (Reviewer R2.9), we incorporate a noise profile calibrated to the \textit{IBM Heron r2} processor (e.g., ibm_torino). The noise model includes:
\begin{itemize}
    \item \textbf{Thermal Relaxation:} Characteristic times $T_1 \approx 250\mu s$ and $T_2 \approx 150\mu s$.
    \item \textbf{Gate Depolarizing Error:} $p_{1q} = 2 \cdot 10^{-4}$ for single-qubit gates and $p_{2q} = 5 \cdot 10^{-3}$ for CNOTs.
    \item \textbf{Readout Error:} Measurement assignment errors of $1.2\%$.
\end{itemize}

\subsection{Statistical Protocol and Metrics}
Addressing the lack of statistical significance (Reviewer R3.1), every experimental configuration is executed across \textbf{5 independent random seeds}. Results are reported as $\mu \pm \sigma$ (Mean $\pm$ Standard Deviation). Beyond raw Accuracy, we evaluate:
\begin{itemize}
    \item \textbf{F1-Score and ROC-AUC:} For medical dataset robustness (Brain Tumor, Hymenoptera).
    \item \textbf{Energy Consumption (kWh):} Measured via CodeCarbon for environmental and computational efficiency analysis.
\end{itemize}
